\documentclass[11pt]{article}
\usepackage[a4paper,margin=27mm]{geometry}
\usepackage{amsmath,amssymb,amsthm,array,booktabs}
\usepackage[T1]{fontenc}
\usepackage{lmodern}
\usepackage[colorlinks=true,linkcolor=blue,urlcolor=blue,citecolor=blue]{hyperref}
\setlength{\emergencystretch}{1em}
\newtheorem{theorem}{Theorem}
\newtheorem{lemma}[theorem]{Lemma}
\newtheorem{corollary}[theorem]{Corollary}
\newcommand{\Z}{\mathbb Z}
\newcommand{\N}{\mathbb N}
\newcommand{\diam}{\operatorname{diam}}
\title{Extending Meddaugh and Stephens' examples to every finite neighborhood-shadowing separation}
\author{Haobo Ma\\\small ChronoAI Pte Ltd, Singapore
\and Wenlin Zhang\\\small National University of Singapore\\\small ChronoAI Pte Ltd, Singapore}
\date{November 19, 2026}
\begin{document}
\maketitle
\begin{abstract}
Extending the $N=1$ examples of Meddaugh and Stephens,
\emph{Shadowing with small diameter sets of bounded cardinality}
(arXiv:2509.04312v1, September 4, 2025, Examples 3.4 and 3.6,
Theorem 3.7), we answer their Questions 5.1 and 5.2
(Section 5, printed page 11) for every $N\geq1$ by mixing binary
subshifts with dense periodic points.
For each $m\geq3$, we construct a compact two-sided subshift with
neighborhood $(m-1)$-shadowing but without neighborhood $(m-2)$-shadowing.
A periodic phase cover supplies one fixed tracing set of at most $m-1$
points with arbitrarily small initial diameter. A progression with widely
separated occupied sites forces the matching lower bound.
Every member fails quasi-finite type and is nonsofic; in particular,
$m=3$ gives another affirmative answer to Question 5.3.
Moreover, the tracing points can all follow the pseudo-orbit at common
periodic times, so even uniform syndetic tracing by each point does not
characterize quasi-finite type.
\end{abstract}

\section{The question and the construction}
The questions of Meddaugh and Stephens~\cite[Questions 5.1--5.2]{MS}
ask which strict finite levels of neighborhood shadowing occur in dynamical
systems and in subshifts. All positive finite levels occur, already in
binary two-sided subshifts. Their positive results for mixing sofic shifts
and shifts of quasi-finite type give neighborhood $2$-shadowing
\cite[Theorems 3.2 and 3.7]{MS}; the construction here allows any prescribed
finite bound, including $3$-shadowing without $2$-shadowing.
The separation between $2$-shadowing and ordinary shadowing was already
obtained in their Examples 3.4 and 3.6 and Theorem 3.7.
We also prove that every member is mixing and has dense periodic points,
although none is sofic.

We use exactly the initial-diameter condition in
\cite[Definition 2.1]{MS}. A continuous map $f$ on a compact metric space
has \emph{neighborhood $r$-shadowing} if for every $\varepsilon>0$ there
is $\delta>0$ such that each sequence $(z_t)_{t\geq0}$ satisfying
$d(f(z_t),z_{t+1})<\delta$ admits a nonempty finite set $A$ with
\[
 |A|\leq r,\qquad \diam A<\varepsilon,\qquad
 \min_{a\in A}d(f^t(a),z_t)<\varepsilon\quad(t\geq0).
\]
The set $A$ is fixed; its matching member may depend on $t$.
Only its initial diameter is constrained.

For $m\geq3$ put $b=m+1$. For $s\in\Z$ and $k\in\N=\{0,1,\ldots\}$,
let
\[
 \begin{gathered}
 P_{s,k}=\{s+j b^k:0\leq j<m\},\\
 X_m=\left\{x\in\{0,1\}^{\Z}:
       \sum_{j=0}^{m-1}x_{s+j b^k}\leq2
       \text{ for all }s\in\Z,\ k\in\N\right\}.
 \end{gathered}
\]
Write $(\sigma x)_n=x_{n+1}$. On the full shift use the metric
$d(x,x)=0$ and, for $x\ne y$,
\[
 d(x,y)=2^{-\min\{|n|:x_n\ne y_n\}}.
\]
Agreement on $[-R,R]$ gives distance at most $2^{-(R+1)}$;
distance less than $2^{-R}$ forces agreement there.
The space $X_m$ contains the zero configuration, is closed because each
constraint involves finitely many coordinates, and is invariant under
every integer shift. Hence it is compact and $\sigma|_{X_m}$ is a
homeomorphism, with inverse the shift by $-1$.

\begin{theorem}\label{thm:main}
For every $m\geq3$, $(X_m,\sigma)$ has neighborhood $(m-1)$-shadowing
and does not have neighborhood $(m-2)$-shadowing. Consequently every
$N\geq1$ occurs in both Questions 5.1 and 5.2, by taking $m=N+2$.
\end{theorem}

\section{A phase cover with a common initial neighborhood}
An interval below denotes its open image in the circle $\mathbb R/m\Z$.
For $m=3$, take
\[
 B_0=(-1/4,7/4),\qquad B_1=(m-7/4,m+1/4).
\]
For $m\geq4$ retain these two sets, put $h=1-1/(2(m-3))$, and for
$2\leq a\leq m-2$ put
\[
 c_a=7/4+(a-3/2)h,\qquad
 B_a=(-1/4,1/4)\ \cup\ (c_a-1/2,c_a+1/2).
\]

\begin{lemma}\label{lem:circle}
The $m-1$ open sets $B_a$ cover $\mathbb R/m\Z$, all contain
$(-1/4,1/4)$, and each contains at most two points of
$\{u,u+1,\ldots,u+m-1\}$ for every real $u$.
\end{lemma}
\begin{proof}
An open interval of length at most $1$ contains at most one of these
points, and one of length at most $2$ contains at most two. Indeed,
their lifts differ by distinct integers, so three lifts span at least
$2$. This proves the count for $B_0,B_1$ and for each union defining
$B_a$. The common neighborhood is immediate, using the lift around $m$
for $B_1$.

If $m=3$, the two length-$2$ arcs overlap and cover the circle.
For $m\geq4$, the remote intervals have successive centers separated by
$h<1$, so their union is one open interval. Its left endpoint is
$5/4+h/2<7/4$, and its right endpoint is
\[
 9/4+(m-7/2)h=m-7/4+\frac1{4(m-3)}>m-7/4.
\]
Thus they bridge the gap between $B_0$ and $B_1$. These two arcs also
overlap across $0$ modulo $m$, proving the cover.
\end{proof}

\begin{lemma}\label{lem:masks}
For every integer $R\geq0$ there are $K\geq0$ and sets
$M_0,\ldots,M_{m-2}\subseteq\Z$ such that:
\begin{enumerate}
\item every $M_a$ contains $[-R,R]\cap\Z$;
\item every window $[t-R,t+R]\cap\Z$, $t\geq0$, lies in some $M_a$;
\item $|M_a\cap P_{s,k}|\leq2$ for all $s\in\Z$ and $k\geq K$.
\end{enumerate}
\end{lemma}
\begin{proof}
The finite open cover in Lemma~\ref{lem:circle} has a uniform window
margin: there is $\eta>0$ such that each circle point has its
$\eta$-neighborhood contained in some $B_a$. To see this without a
choice depending on infinitely many windows, choose for each point a
ball whose double-radius ball lies in a cover member. A finite subcover
of the smaller balls exists by compactness; the minimum of their radii
works. Decrease $\eta$ so that $\eta<1/4$.
Choose $K$ with $Q=b^K>R/\eta$, and set
\[
 M_a=\{n\in\Z:n/Q\bmod m\in B_a\}.
\]
The first two claims follow from $R/Q<\eta$ and the common
neighborhood at $0$. For $k\geq K$, $b^{k-K}\equiv1\pmod m$.
The phases of the vertices of $P_{s,k}$ are therefore
$s/Q+j\pmod m$, $0\leq j<m$. Lemma~\ref{lem:circle} proves the last
claim, including every translated progression.
\end{proof}

\section{The upper bound}
Fix $\varepsilon>0$ and choose $R\geq0$ with
$2^{-(R+1)}<\varepsilon$. Take the masks and $K$ from
Lemma~\ref{lem:masks}. Put $H=(m-1)b^K$; enlarging $K$ if necessary
ensures $H\geq R$. Let $\delta=2^{-H}$ and let $(z_t)_{t\geq0}$ be
a $\delta$-pseudo-orbit in $X_m$. Its adjacent states satisfy
\begin{equation}\label{eq:edge}
 (z_t)_{j+1}=(z_{t+1})_j\qquad(-H\leq j\leq H).
\end{equation}
Define a single signed configuration by
\[
 c_n=\begin{cases}(z_0)_n,&n<0,\\(z_n)_0,&n\geq0.\end{cases}
\]
Repeated use of \eqref{eq:edge} gives
\begin{equation}\label{eq:overlap}
 c_{t+j}=(z_t)_j\quad(t\geq0,\ |j|\leq H),
 \qquad c_n=(z_0)_n\quad(n\leq H).
\end{equation}
For $j\geq0$, propagate from time $t$ to $t+j$ with decreasing
coordinate. For $j<0$ and $t+j\geq0$, propagate backwards to time
$t+j$ with increasing coordinate. If $t+j<0$, propagate backwards
to time $0$; every coordinate used lies between $j$ and $0$.
These are precisely the ranges permitted by \eqref{eq:edge}.

Every progression constraint of $c$ at scale $k<K$ is legal.
Indeed, its span $D=(m-1)b^k$ is at most $H$. If its largest
coordinate is at most $H$, the whole progression belongs to the
unchanged part of $z_0$ in \eqref{eq:overlap}. Otherwise its smallest
coordinate $s$ is positive, since $s+D>H$ and $D\leq H$.
Its vertices agree with coordinates $0,b^k,\ldots,(m-1)b^k$ of
$z_s$, by \eqref{eq:overlap}. In either case legality follows from
membership of that state in $X_m$.

For $0\leq a\leq m-2$, define an entire configuration
\[
 y^a_n=\begin{cases}c_n,&n\in M_a,\\0,&n\notin M_a.\end{cases}
\]
Deleting ones cannot violate a short-scale constraint. At scale $k\geq K$,
at most two vertices belong to $M_a$, by Lemma~\ref{lem:masks}.
Thus each $y^a$ belongs to $X_m$. Take
$A=\{y^0,\ldots,y^{m-2}\}$, with repeated configurations counted once.
All its members agree with $z_0$ on $[-R,R]$, so
$\diam A\leq2^{-(R+1)}<\varepsilon$. For every $t\geq0$,
some mask contains $[t-R,t+R]$; for that $a$,
\[
 (\sigma^t y^a)_j=c_{t+j}=(z_t)_j\quad(|j|\leq R).
\]
Hence $d(\sigma^t y^a,z_t)\leq2^{-(R+1)}<\varepsilon$.
This proves neighborhood $(m-1)$-shadowing with one fixed nonempty set.

\section{The matching lower bound}
Fix $\varepsilon_0=1/2$. Given any $\delta>0$, choose $L\geq1$
with $2^{-L}<\delta$ and choose $k$ with $q=b^k>2L+2$.
Let $V=\{0,q,\ldots,(m-1)q\}$ and define $z_t$ to be the
indicator configuration of
\[
 (V-t)\cap[-L,L]\quad(t\geq0).
\]
This support has at most one point, so $z_t\in X_m$.
For $-L+1\leq j\leq L-1$,
\[
 (\sigma z_t)_j=1\ \Longleftrightarrow\ t+j+1\in V
 \ \Longleftrightarrow\ (z_{t+1})_j=1.
\]
Consequently $d(\sigma z_t,z_{t+1})\leq2^{-L}<\delta$.
Also $(z_{jq})_0=1$ for $0\leq j<m$.

Suppose a finite nonempty $A$ with $\diam A<1/2$ traces this sequence
to accuracy $1/2$. At time $0$ some member has coordinate $0$ equal
to $1$. The diameter condition forces every member of $A$ to have
that same coordinate equal to $1$. For each $1\leq j<m$, tracing
at time $jq$ requires some $a\in A$ with $a_{jq}=1$.
No single $a$ can serve two distinct such indices: together with
$a_0=1$ they would give three ones on $P_{0,k}$, contrary to
$a\in X_m$. Thus $m-1$ different members are necessary. In particular
no set of at most $m-2$ members works. This proves the lower bound
and completes Theorem~\ref{thm:main}.

\section{Consequences and coverage}
\begin{corollary}\label{cor:rank}
The subshift $X_4$ has neighborhood $3$-shadowing but not neighborhood
$2$-shadowing. More generally the least finite neighborhood-shadowing
bound of $X_m$ is exactly $m-1$.
\end{corollary}
\begin{proof}
Apply Theorem~\ref{thm:main}. A set of cardinality at most $r$ is also
one of cardinality at most $r+1$, so failure at $m-2$ excludes every
smaller bound.
\end{proof}

\begin{corollary}\label{cor:qft}
No $X_m$ is of quasi-finite type, or of finite type. In particular,
$X_3$ answers Question 5.3 of~\cite{MS} affirmatively.
\end{corollary}
\begin{proof}
Recall \cite[Definition 3.5]{MS}: quasi-finite type requires a bound $M$
such that, whenever words $u,w,v$ belong to the language with $uw,wv$
allowed and $|w|\geq M$, some word $z$ with $|z|=|w|$ makes $uzv$
allowed. Given any proposed $M$, choose $q=b^k$ with $q-1\geq M$,
and put
\[
 u=10^{q-1}1,\qquad w=0^{q-1},\qquad v=1.
\]
Both $uw$ and $wv$ occur in configurations with at most two ones,
so they belong to the language of $X_m$. But every $uzv$ with
$|z|=q-1$ has ones at positions $0,q,2q$. Those are three vertices of
$P_{0,k}$ since $m\geq3$, so no such word belongs to the language.
This disproves quasi-finite type. A shift of finite type is of
quasi-finite type: if the longest forbidden word has length $B$ and
$|w|\geq B$, any forbidden word in $uwv$ would be contained in $uw$
or in $wv$. Splicing a left extension of $uw$ and a right extension of
$wv$ through the shared word $w$ gives a full configuration avoiding
all forbidden words. Thus $uwv$ is allowed, and $X_m$ is not of finite type.
Theorem~\ref{thm:main} supplies neighborhood $2$-shadowing for $X_3$.
\end{proof}

The companion note of Ma and Zhang~\cite{MZ} answers Question 5.3
by a distinct return-distance construction. Its specialization to powers
of two is stated by \nolinkurl{powers_two_return_distance_example}
in the pinned
\href{https://github.com/the-omega-institute/trureturing/blob/147008bff5524a3dccf0efa3ed9e76fb6b168a50/D5/S1/Dynamics/Shadowing/PowersTwoReturnDistances.lean}{\texttt{PowersTwoReturnDistances.lean}}.
Corollary~\ref{cor:qft} instead concerns the progression family $X_m$.
The powers-of-two example allows the all-ones configuration, whereas
$X_3$ forbids $111$.

\begin{corollary}\label{cor:returns}
For every $m\geq3$ and every $\varepsilon>0$, there are $\delta>0$
and an integer $p\geq1$ such that every $\delta$-pseudo-orbit $(z_t)$
in $X_m$ admits a fixed nonempty set $A$ with $|A|\leq m-1$,
$\diam A<\varepsilon$, and
\[
 \min_{a\in A}d(\sigma^t a,z_t)<\varepsilon\quad(t\geq0),
 \qquad
 \max_{a\in A}d(\sigma^{\ell p}a,z_{\ell p})<\varepsilon
 \quad(\ell\geq0).
\]
In particular, even neighborhood two-shadowing in which each tracer
has a uniformly bounded return time to the pseudo-orbit does not imply
quasi-finite type.
\end{corollary}
\begin{proof}
The masks in Lemma~\ref{lem:masks} are periodic with common integer
period $p=mQ$, since adding $mQ$ to a coordinate adds $m$ to its phase.
Each mask contains $[-R,R]$, so each contains
$[\ell p-R,\ell p+R]$ for every integer $\ell$.
The upper-bound construction therefore makes every $y^a$ agree with
$z_{\ell p}$ on the radius-$R$ observation window at time $\ell p$.
The choices of $p$ and $\delta$ depend on $\varepsilon,m$, not on
the pseudo-orbit. Each block of $p$ consecutive natural times contains
a multiple of $p$, giving the asserted uniform return bound.
For the last claim take $m=3$ and use Corollary~\ref{cor:qft}.
\end{proof}

The discussion immediately before Question 5.3 of~\cite{MS} notes
syndetic tracing in the quasi-finite-type construction. The converse
with a uniform syndetic bound is not asserted there;
Corollary~\ref{cor:returns} shows that this natural strengthening of
two-shadowing alone is insufficient for a characterization.

\begin{corollary}\label{cor:mixing}
For every $m\geq3$, $X_m$ is topologically mixing and has dense periodic
points. It is nonsofic. Thus every separation in Theorem~\ref{thm:main}
occurs within mixing nonsofic subshifts with dense periodic points.
\end{corollary}
\begin{proof}
Deleting ones preserves every defining constraint. Hence the zero
extension of any allowed word is legal. Let $u,v$ be allowed nonempty
words, put $D=\max(|u|-1,|v|-1)$, and separate their zero extensions by
a zero gap whose length exceeds $(m-1)D$. If a progression contained
three ones in their union, it would meet both word intervals and two of
its occupied vertices would lie in one interval. Its step $b^k$ would
then be at most $D$, so its total span would be at most $(m-1)D$.
This is smaller than the distance between the intervals, a contradiction.
The union is legal for every sufficiently long gap, which proves mixing
on cylinders and therefore on arbitrary nonempty open sets.

For density of periodic points, let $u$ be an allowed word of length
$L$, choose $Q=b^K>L$, and repeat $u$ with zeros between its copies at
period $p=mQ$. Call the resulting point $x$. A progression of scale
$k<K$ has span less than $(m-1)Q$, whereas different word intervals are
separated by distance at least $mQ-(L-1)>(m-1)Q$. It meets at most one
copy of $u$, where the zero extension is legal. For $k\geq K$, its step
is congruent to $Q$ modulo $mQ$, since $b\equiv1\pmod m$. Its $m$
vertices modulo $p$ form a full set spaced by $Q$. An interval of length
$L<Q$ contains at most one of those vertices. Hence these progressions
also contain at most two ones, and $x\in X_m$. Every allowed cylinder
therefore contains a periodic point. Its singleton is a nonempty compact
periodic set, so $X_m$ has dense small periodic sets as well.

Finally, a mixing sofic shift has a uniform mixing number
\cite[Proposition 2.3]{MS}: for all allowed $u,v$ and every sufficiently
large length $n$, a word $z$ of length $n$ makes $uzv$ allowed.
This implies quasi-finite type, contradicting Corollary~\ref{cor:qft}.
Thus $X_m$ is not sofic.
\end{proof}

\noindent The following table refers to Section 5 of~\cite{MS}.
\begin{center}
\begin{tabular}{@{}>{\raggedright\arraybackslash}p{.22\linewidth}>{\raggedright\arraybackslash}p{.71\linewidth}@{}}
\toprule
Source item&Coverage\\\midrule
Questions 5.1, 5.2 (p.\ 11)&Answered: every $N\geq1$, by $X_{N+2}$.
For $N=1$, reproved; prior examples of Meddaugh--Stephens (2025).
All witnesses here are mixing, nonsofic, with dense periodic points.\\
Preceding $3$ versus $2$ case (p.\ 11)&Answered: yes, by $X_4$.\\
Question 5.3 (p.\ 11)&Answered: yes, by $X_3$; the companion
return-distance example is in~\cite{MZ}.\\
Preceding syndetic-tracing comment&Not asserted in the source: the
uniform-syndetic converse fails for $X_3$, even with common periodic
tracing times (Corollary~\ref{cor:returns}).\\
Questions 5.4, 5.5&Not covered: an explicit family is not a characterization
of either class of subshifts.\\
Question 5.6 (p.\ 12)&Not covered by this family: its periodic points
are dense. A separate example answering yes is specified below.\\\bottomrule
\end{tabular}
\end{center}

Question 5.6 has a separate affirmative construction at $N=2$, stated by
\nolinkurl{quadratic_progression_example} in
\href{https://github.com/the-omega-institute/trureturing/blob/147008bff5524a3dccf0efa3ed9e76fb6b168a50/D5/S1/Dynamics/Shadowing/ProgressionExample.lean}{\texttt{ProgressionExample.lean}}.
It gives a compact surjective weakly mixing system with neighborhood
$2$-shadowing and a nonempty clopen set containing no nonempty set
mapped into itself by any positive forward power. In particular, that open set
contains no nonempty compact periodic set. Corollary~\ref{cor:mixing}
explains why the family $X_m$ is different.

\paragraph{Formal statement and finite checks.}
The pinned Lean declaration
\nolinkurl{lacunary_shadowing_hierarchy} in
\href{https://github.com/the-omega-institute/trureturing/blob/147008bff5524a3dccf0efa3ed9e76fb6b168a50/D5/S1/Dynamics/Shadowing/LacunaryShadowingHierarchy.lean}{the Lean companion}
states the nonemptiness, compactness, continuity, bijectivity and both
neighborhood-shadowing conclusions of Theorem~\ref{thm:main}.
Its predicate uses a fixed nonempty finite tracing set and initial diameter,
with natural-time pseudo-orbits on integer-coordinate configurations.
The supporting declarations have the following separate scopes:
\begin{itemize}
\item In the \href{https://github.com/the-omega-institute/trureturing/blob/147008bff5524a3dccf0efa3ed9e76fb6b168a50/D5/S1/Dynamics/Shadowing/LacunaryProgressionSubshift.lean}{progression-subshift module},
\nolinkurl{NeighborhoodFiniteShadowing} and \nolinkurl{LacunaryLegal}
define the tracing property and the progression constraint;
\nolinkurl{lacunary_legal_closed}, \nolinkurl{lacunary_legal_shift},
\nolinkurl{lacunary_zero_legal} and \nolinkurl{lacunary_shift_continuous}
give closure, translation invariance, nonemptiness and continuity.
\item In \href{https://github.com/the-omega-institute/trureturing/blob/147008bff5524a3dccf0efa3ed9e76fb6b168a50/D5/S1/Dynamics/Shadowing/LacunaryPhaseCover.lean}{\texttt{LacunaryPhaseCover.lean}},
\nolinkurl{lacunary_phase_cover} supplies $m-1$ masks containing the
initial window, covering every later window, and avoiding triples on
every translated sufficiently large progression.
\item In \href{https://github.com/the-omega-institute/trureturing/blob/147008bff5524a3dccf0efa3ed9e76fb6b168a50/D5/S1/Dynamics/Shadowing/LacunaryShadowingUpper.lean}{\texttt{LacunaryShadowingUpper.lean}},
\nolinkurl{lacunary_neighborhood_upper} proves neighborhood
$(m-1)$-shadowing for every $m\geq3$.
\item In \href{https://github.com/the-omega-institute/trureturing/blob/147008bff5524a3dccf0efa3ed9e76fb6b168a50/D5/S1/Dynamics/Shadowing/LacunaryShadowingLower.lean}{\texttt{LacunaryShadowingLower.lean}},
\nolinkurl{lacunary_not_neighborhood_smaller} excludes neighborhood
$(m-2)$-shadowing for every $m\geq3$.
\item In \href{https://github.com/the-omega-institute/trureturing/blob/147008bff5524a3dccf0efa3ed9e76fb6b168a50/D5/S1/Dynamics/Shadowing/ReturnDistanceShadowing.lean}{\texttt{ReturnDistanceShadowing.lean}},
\nolinkurl{return_metric_window} and \nolinkurl{return_window_control}
relate finite coordinate windows to the product metric, and
\nolinkurl{signed_diagonal_overlap} proves the signed diagonal identity.
\end{itemize}
The formal metric is the compatible countable-product metric; the metric
used above is uniformly equivalent on the compact space. The property is
unchanged by such a replacement: transfer the requested accuracy by uniform
continuity in one direction, obtain a tracing tolerance, then transfer that
tolerance by uniform continuity in the other direction. This controls
diameters and pointwise tracing by the same two implications.
The supplied exact-rational checks cover the circle masks for
$3\leq m\leq20$. The arguments above prove the claims for every $m\geq3$;
Corollaries~\ref{cor:qft}, \ref{cor:returns} and~\ref{cor:mixing}
have written proofs here.
The independent integer checks in \texttt{verify-structure.py} cover all
allowed binary words of length at most $6$ for $3\leq m\leq6$: they
test zero-gap gluing, periodic extensions, and the exact-length bridge
obstruction. Their finite ranges supplement the all-parameter proofs.

\begin{thebibliography}{9}
\bibitem{MS} Jonathan Meddaugh and Elyssa Stephens,
\emph{Shadowing with small diameter sets of bounded cardinality},
arXiv:2509.04312v1, September 4, 2025.
\href{https://arxiv.org/abs/2509.04312v1}{arXiv:2509.04312v1}.
\bibitem{MZ} Haobo Ma and Wenlin Zhang,
\emph{Neighborhood two-shadowing without quasi-finite type},
unpublished mathematical note, October 2026, Theorem 1.
The powers-of-two statement is given in
\href{https://github.com/the-omega-institute/trureturing/blob/147008bff5524a3dccf0efa3ed9e76fb6b168a50/D5/S1/Dynamics/Shadowing/PowersTwoReturnDistances.lean}{\texttt{PowersTwoReturnDistances.lean}}.
\end{thebibliography}
\smallskip\noindent\small
\copyright\ 2026 Haobo Ma and Wenlin Zhang.
Licensed under Creative Commons Attribution 4.0 International (CC BY 4.0).
\end{document}
